% THETA3-B25 contribution by /root/definite_bridge, 21 November 2026.
\subsection{THETA3-B25: the actual theta image has an algebraic finite-period inverse}\label{B25}

The finite-source argument of B24 has a part that does not
require the source to be totally definite.  Finite adelic
source translates already span the actual finite theta
image.  Their common CM evaluation scalar proves that
this image is algebraic in the coherent coordinates of
B23.  Consequently the algebraic finite-period correction
can be chosen to preserve the original theta image.
The integral source Gram matrix of B17 is not replaced
by this finite restriction matrix.

\paragraph{Fixed kernel and the finite-adelic spanning lemma.}
Retain all the data of B17, including the full adelic
target realization, its scalar canonical type, the
coefficient pairing, additive character, splittings,
finite levels and Haar measures.  In particular
\[
 K(g,h)=\theta_{V,P}(g,h;\Phi),\qquad
 Tu=\int_{[S]}K(\,\cdot\,,h)u(h)\,dh,
 \qquad \mathcal E_0=\operatorname{im}T,
 \quad S=\U(P).
 \tag{B25a}
\]
The source $P$ is the actual rank-two multiplicity
plane, definite at the active place but possibly
indefinite elsewhere.  Both quotients are compact,
and B17 proves that the image is finite dimensional
and is the span of all $K(\,\cdot\,,h)$.
Assume, as in B24, that the fixed finite input has
algebraic coefficients and that the finite splitting
values used in the finitely many translations below
are algebraic.  All specified original complex
normalization scalars are retained.

Then
\[
 \mathcal E_0=
  \operatorname{span}_{\C}
       \{K(\,\cdot\,,(1,b)):b\in S(\A_{F,f})\}.
 \tag{B25b}
\]
In particular, if $r=\dim\mathcal E_0>0$, one can
choose $r$ such finite adelic elements whose kernel
vectors are a basis.

\textbf{Proof.}
The Cayley transform identifies open sets of $S$
and its $F$-vector space of skew-adjoint matrices:
$X\mapsto(1-X)(1+X)^{-1}$, with the same formula
for its inverse.  Weak approximation for the
coefficients of that vector space shows that
$S(F)$ is dense in $S(F\otimes_\Q\R)$; the omitted
real locus having eigenvalue $-1$ is a proper
algebraic subset and has empty interior.  This is
an assertion for the full unitary group, not an
assumption about a particular special-unitary
arithmetic subgroup.

Let $h=(h_\infty,h_f)$.  Choose
$\gamma_j\in S(F)$ with
$\gamma_{j,\infty}\to h_\infty$.  Rational theta
automorphy gives the exact vector equality
\[
 K(\,\cdot\,,(1,\gamma_{j,f}^{-1}h_f))
       =K(\,\cdot\,,(\gamma_{j,\infty},h_f)).
 \tag{B25c}
\]
The right side converges to $K(\,\cdot\,,h)$ in
the finite-dimensional holomorphic target space,
by the smooth coefficient dependence proved in A7.
The span on the right of (B25b) is a linear
subspace of a finite-dimensional vector space,
so it is closed.  It therefore contains every
kernel vector.  The reverse inclusion is immediate.
Selecting a basis proves the last assertion.
No uniform-continuity assertion over varying finite
representatives is needed in this argument: $h_f$
is fixed on the right of (B25c). \hfill$\square$

\paragraph{The one CM type and the common finite-value scalar.}
Fix the original rational decomposition
$P=P_1\perp P_2$, with complete types $v_1,v_2$.
Each of the finitely many vectors in (B25b) has the
same archimedean coefficient as the original kernel.
Its finite input is the corresponding finite Weil
translate of $\Phi_f$.  The finite generator
calculation in B24 proves that its coefficients
are algebraic.  Every such input is a finite
algebraic sum of product cosets for $P_1\perp P_2$.
The exact double alternation in 032,
Proposition 4.6(i), printed pp.22--23,
then expresses the kernel vector as a finite sum
of wedges of the two prescribed types.

Apply 032, Corollary 6.7, printed p.45, to this
finite collection of actual inputs.  After one
common principal cover, their degree-one classes
come from finitely many rational maps from the
specified CM modules.  Use B23's algebraic CM
cup source on that cover.  This does not require
that all of $H^1$ be generated by theta forms.

The relevant source tensor is the equal-label
summand $U(v_1)\otimes_E U(v_2)$, obtained by the
diagonal idempotent in $E\otimes_\Q E$.
It has rank one over $E$.  At an embedding its
Hodge type is determined by $v_1+v_2$ and is one
of $(2,0),(1,1),(0,2)$.  Its Mumford--Tate
characters form one Galois orbit, possibly with
repetitions, since the embeddings of $E$ form
one orbit.  It is therefore isotypic as a
rational CM Hodge structure.  The active
selected line has type $(2,0)$, so its simple
type is non-Tate.  Taking rational source maps
and then cup product preserves this assertion.
Consequently all the selected kernel vectors,
and hence all of $\mathcal E_0$, lie in one
full multiplicity block $\mathcal E_\chi$ of B23.

In the coherent basis $\omega=P_\chi\beta$ of
that block, choose A5's separating finite
target evaluation rows.  Their matrix is
$t_xA$, with $A$ algebraic and of full column
rank.  Evaluate the finitely many source
basis vectors at these rows.  The archimedean
positive Lagrangian and rational CM base point
are fixed; only finitely many algebraic finite
Weil inputs have changed.  B19 therefore gives
one nonzero comparison scalar $C_x$, with its
same theta-module, tangent, level and constant-line
factors, such that their evaluation matrix is
$C_xB$, $B\in M(\overline\Q)$.

Let $E$ be the coefficient matrix of these
kernel vectors in the basis $\omega$.  For
any algebraic left inverse $L_A$ of $A$,
\[
       E=\frac{C_x}{t_x}L_AB.
 \tag{B25d}
\]
Thus $\mathcal E_0$ is a subspace defined over
$\overline\Q$ in coherent coordinates.  Choose
an algebraic full-rank column matrix $F$ whose
columns are a basis of that subspace.
This proof uses actual CM-point theta values;
it does not treat evaluation as a morphism of
weight-two Hodge structures.  Nor does it
assert that the original, unnormalized theta
basis vectors themselves have algebraic
coherent coordinates: their common factor
$C_x/t_x$ remains in (B25d).

\paragraph{A correction preserving the original image.}
Write
\[
 G_0=F^*G_BF,\qquad A_0=AF,\qquad
 H_0=A_0^*DA_0,
 \tag{B25e}
\]
where $D$ is the finite positive diagonal
probability-mass matrix for the definite
restrictions.  It has rational entries, with
the original masses retained if stabilizer
conditions force some rows to vanish.
Both $G_0$ and $H_0$ are positive definite
algebraic Hermitian matrices.
Let $Q_{\rm fin}$ be A5's actual restriction
operator on $\mathcal E_0$, with its
inherited Petersson pairing.  In the basis
$\omega F$ it satisfies
\[
 \begin{split}
 Q_{\rm fin}&=
  \frac{|t_x|^2}{\nu_G|P_\chi|^2}G_0^{-1}H_0,\\
 Q_{\rm fin}^{-1}&=s_\chi K_0,
 \qquad s_\chi=\frac{\nu_G|P_\chi|^2}{|t_x|^2},
 \qquad K_0=H_0^{-1}G_0.
 \end{split}
 \tag{B25f}
\]
These identities follow by taking the adjoint
with domain Gram matrix $\nu_G|P_\chi|^2G_0$
and range Gram matrix $D$.

The algebraic matrix $K_0$ extends by zero
on the $G_B$-orthogonal complement of the
column space of $F$.  That projection has
matrix $F(F^*G_BF)^{-1}F^*G_B$.
B23's CM-source retract realizes the extended
matrix by an algebraic correspondence with
$\overline\Q$ coefficients on the target.
In particular $K_0$ preserves this very
$\mathcal E_0$, not merely its full containing
CM block.

Choose the definite test functions and
periods $\mathcal P_{j,\alpha}$ exactly as
in A5.  Then, for all $f,g\in\mathcal E_0$,
\[
 \langle f,g\rangle_G
   =s_\chi\sum_{j,\alpha}
       \overline{\mathcal P_{j,\alpha}(K_0f)}
                        \mathcal P_{j,\alpha}(g).
 \tag{B25g}
\]
The tests are actual totally definite rank-two
subgroup periods with the stated Haar measures.
Every corrected form $K_0f$ still lies in the
image of the original fixed theta operator.
Thus B17's actual adjoint double lift and
literal last-step seesaw apply to it using
the same $P,V,\Phi$ and splittings.
This closes the image-preservation condition
in this finite-period comparison.  The
single $K_0$ and $s_\chi$ work for all A7
source characters with these fixed data.

\paragraph{The actual fixed-data variation in 032.}
The parameter space in 032, Section 4.6,
printed pp.22--24, varies only the positive
and negative lines of the same multiplicity
plane at the compact target places.  Its
active positive Lagrangian, rational total
space, finite representation and total finite
input are fixed.  Each such line space is
an orbit of the appropriate $\U(P_\tau)$.
Choose $s_t\in S(F\otimes_\Q\R)$ carrying
the initial lines to those at $t$, and choose
its active component to be the identity.
In a parameter neighborhood it can be
chosen smoothly.

Heisenberg covariance transports the
annihilator equations for the Gaussian.
Consequently the specified Weil action
$\omega(s_t)$ carries the initial nonzero
Gaussian vacuum line to the new Gaussian
vacuum line.  Their chosen vectors differ
by a nonzero scalar.  The active determinant
cochain is unchanged because $s_{t,\mu}=1$.
The selected determinant half-lines and
their square-root frames add only their
recorded nonzero scalar.  These factors
are independent of the target ball variable:
only the compact-place data have changed,
and the active data are identical.  Equivalently
the two coefficient identifications have an
invariant ratio under the transitive target
special-unitary action, hence a constant
ratio, exactly as in the proof of 032,
Proposition 4.6(i).
Thus, with the original splitting retained,
the literal scalar coefficient satisfies
\[
       \Phi_t=c_t\omega(s_t)\Phi_0,
       \qquad
       \Xi(\Phi_f;t)=c_tK(\,\cdot\,,s_t),
       \qquad c_t\ne0.
 \tag{B25h}
\]
Here $c_t$ is not asserted algebraic or
silently set equal to one.  This statement
uses A7's common closed-cochain realization;
the continuity assertion of 032 alone did
not claim holomorphy at intermediate $t$.
It follows from (B25h) and A7 that every
such form belongs to the same $\mathcal E_0$.

There is also an independent check of the
intermediate holomorphy using only 032's
displayed conclusions.  Rational nonisotropic
vectors of $P(E)$ can approximate prescribed
nonisotropic vectors at all real places
simultaneously.  Their lines and orthogonal
complements give rational decompositions
dense in this parameter product; the active
definite place imposes no additional sign
condition.  Proposition 4.6(i) makes the
corresponding forms holomorphic, and (iii)
gives convergence in $C^\infty$ on the fixed
quotient.  Passing to the limit in
$\bar\partial\Xi_t$ proves it is zero.
An individual wedge expansion can be checked
on its deeper principal cover; vanishing of
$\bar\partial$ then descends to the fixed
quotient.  This explains why no holomorphy
assumption at intermediate parameters is
being added.

In particular the rational balanced
endpoints of 032, Corollary 4.7, printed
pp.24--25, use the same preserved-image
correction (B25f).  Their source tori may
be different rational subgroups of this
same $S$; integrating any of the resulting
torus functionals still takes values in
the same finite image.
Balance also identifies their complete
rank-one $E$ tensor Hodge types, since it
is exactly equality of the functions
$v_1+v_2=v_3+v_4$.

\paragraph{The remaining integral Gram.}
It is essential to distinguish (B25f) from
the inverse of $Q_T=TT^*$ in B17.
The latter still involves the actual matrix
$W_{ab}=\int_{[S]}e_a\overline{e_b}\,dh$.
The present result makes its coefficient
space algebraic, but does not prove that
this integral matrix is projectively algebraic.
In a last-step application to $K_0f$, the
reverse input $Q_T^{-1}K_0f$ is still the
analytic inverse specified in B17.
Nor is a finite sum of these period identities
one prescribed arithmetic \texttt{MainX}
tuple.  Its exact kernel, lattice and
normalization clauses remain separate.

\paragraph{Framework and source scope.}
The CM correspondence realization is exactly
B23's transpose and inverse-CM-polarization
construction, with the supplied Grothendieck
\emph{cote 44}, scans 7--8, and
\emph{cote 16}, scans 43--46.
The actual analytic source clauses are
\texttt{Development/AbelianVarietyTheory/Definitions/}
\texttt{LatticeFrames/IsQuasiFam.lean},
\texttt{LastLaw}, lines 71--75, and
\texttt{SeeLast}, lines 110--123, as
instantiated with the coefficient pairing
and dagger in B17.  The compact-source
finite sum of B24 is recovered when
$P$ is totally definite; no such finite
source quotient is assumed in (B25b).
